\section{Occupation and incidence moments}
\label{sec:occupation}

Keep the deleted index \(i\), its fixed clause \(C_i\), and the
environment \(\mathcal E\) of Section~\ref{sec:deletion}.  We first
bound the time spent at each level of the shorter-clause killing
probability \(q_m=q_d(S_m)\).  We then combine that occupation bound
with the root tests and average the incidence counts.

In this section \(C_k\) denotes a positive constant depending only on
the fixed \(k\), which may increase between occurrences.  It need not
equal the explicit constant in Proposition~\ref{prop:lifetime}.
Recall \(\beta=k/(k-1)\), \(d=\max\{1,Z\}\), and
\(W_m=\min\{L,\tau_k(S_m)+Z\}\).  On \(S_m\ne\varnothing\),
\begin{equation}\label{eq:weighted-lifetime-product}
 W_m q_m^\beta\le C_kd^\beta.
\end{equation}
Indeed, Proposition~\ref{prop:lifetime} bounds
\(\tau_k(S_m)q_m^\beta\) by a constant times \(d^\beta\), and
\[
 Zq_m^\beta\le Z\le d\le d^\beta.
\]
The product bound remains valid when \(q_m=0\).

\subsection{Time spent above a level}

\begin{lemma}[Conditional occupation]\label{lem:conditional-occupation}
There is a constant \(C_k\ge1\), independent of the environment,
such that for every \(0<a\le1\), almost surely,
\begin{equation}\label{eq:conditional-occupation}
 \sum_{m=i}^{L-1}
 \PP(S_m\ne\varnothing,\ q_m\ge a\mid\mathcal E)
 \le\min\{L,C_kd^\beta a^{-\beta}\}.
\end{equation}
\end{lemma}

\begin{proof}
Define
\[
 N_a=\sum_{m=i}^{L-1}
       \mathbf1_{\{S_m\ne\varnothing,\ q_m\ge a\}},
\]
and let \(\sigma_a\) be the first index counted by \(N_a\), or
\(\infty\) if there is none.  Because \(d\)
is known in \(\mathcal E\) and \(S_m\) is \(\mathcal G_m\)-measurable,
\(\{\sigma_a=m\}\in\mathcal G_m\).

On \(\{\sigma_a=m\}\), no eligible index precedes \(m\).  The future
ordinary clauses used to define \(U_m\) kill the original \(S_m\)
at index \(m+U_m\).  If that index is at most \(L\), then \(S_\ell\)
is empty there and thereafter: a member would have to lie in \(S_m\)
and satisfy those same ordinary clauses.  The eligible indices are
therefore among \(m,\ldots,m+U_m-1\).  If the kill is after \(L\),
the capped interval ends earlier.  Thus in both cases
\[
 N_a\le U_m\quad\text{on }\{\sigma_a=m\}.
\]
The first eligible index is included in these \(U_m\) possible indices.

On \(\{\sigma_a=m\}\), the lifetime bound and
\eqref{eq:chronological-domination} give
\[
 \EE[U_m\mid\mathcal G_m]
 \le\tau_k(S_m)+Z
 \le C_kd^\beta a^{-\beta},
\]
after increasing \(C_k\), since \(q_m\ge a\) and
\(Z\le d\le d^\beta a^{-\beta}\).  The first-visit events are disjoint,
so the tower property yields
\begin{align*}
 \EE[N_a\mid\mathcal E]
 &\le\sum_{m=i}^{L-1}
   \EE\left[\mathbf1_{\{\sigma_a=m\}}
                \EE[U_m\mid\mathcal G_m]\mid\mathcal E\right]\\
 &\le C_kd^\beta a^{-\beta}
       \sum_{m=i}^{L-1}\PP(\sigma_a=m\mid\mathcal E)
 \le C_kd^\beta a^{-\beta}.
\end{align*}
Also \(N_a\le L\).  These two bounds prove the claim.
\end{proof}

Only the first visit to the level is stopped on.  All later eligible
indices are charged to the remaining lifetime of the same nested
set.  This is why the argument needs fresh future ordinary contents
under \(\mathcal G_m\), even though \(\mathcal E\) has already exposed
future nonordinary data.

To apply the occupation bound to the squared deletion delay, first
consider an environment with \(V=0\).  All \(k\) root tests are then
available.  Lemma~\ref{lem:root-flip-tests} supplies the factor
\(q_m^k\) in the weighted-delay bound
\eqref{eq:conditional-weighted-delay}.  On a nonempty \(S_m\),
\eqref{eq:weighted-lifetime-product} gives
\[
 W_mq_m^k=(W_mq_m^\beta)q_m^{k-\beta}
 \le C_kd^\beta q_m^{k-\beta}.
\]
The remaining occupation sum has exponent
\[
 \gamma=k-\beta=\frac{k(k-2)}{k-1}.
\]
The lifetime estimate has controlled the remaining deletion duration;
Lemma~\ref{lem:conditional-occupation} contributes a second lifetime
factor \(a^{-\beta}\) by bounding the time spent above level \(a\).
The resulting layer-cake integral contains \(a^{\gamma-\beta-1}\).
Its integrability at zero is determined by
\(\gamma-\beta=k(k-3)/(k-1)\), which is zero exactly at \(k=3\).

\begin{lemma}[Integrated occupation]\label{lem:integrated-occupation}
For \(L\ge1\), almost surely and with a constant independent of the
environment,
\begin{equation}\label{eq:integrated-occupation}
 \sum_{m=i}^{L-1}
 \EE[\mathbf1_{\{S_m\ne\varnothing\}}q_m^\gamma\mid\mathcal E]
 \le
 \begin{cases}
 C_kd^\beta(1+\log L),&k=3,\\
 C_kd^\beta,&k\ge4.
 \end{cases}
\end{equation}
\end{lemma}

\begin{proof}
Since \(0\le q_m\le1\), the layer-cake identity and
Lemma~\ref{lem:conditional-occupation} give
\begin{align}
 &\sum_{m=i}^{L-1}
 \EE[\mathbf1_{\{S_m\ne\varnothing\}}q_m^\gamma\mid\mathcal E]
 \notag\\
 &\quad=\int_0^1\gamma a^{\gamma-1}
       \sum_{m=i}^{L-1}
       \PP(S_m\ne\varnothing,\ q_m\ge a\mid\mathcal E)\,da
 \notag\\
 &\quad\le\int_0^1\gamma a^{\gamma-1}
        \min\{L,Aa^{-\beta}\}\,da,\qquad
 A=C_kd^\beta\ge1.
 \label{eq:occupation-integral}
\end{align}
The terms are nonnegative, so the interchange of conditional
expectation, the finite sum, and integration is valid.

For \(k\ge4\), \(\gamma>\beta\).  Bounding the minimum by
\(Aa^{-\beta}\) gives
\[
 \int_0^1\gamma a^{\gamma-1}\min\{L,Aa^{-\beta}\}\,da
 \le\frac{\gamma A}{\gamma-\beta}.
\]
This is a constant times \(A\) for fixed \(k\).

For \(k=3\), \(\gamma=\beta=3/2\).  If \(A\ge L\), the integral is
at most \(L\le A\).  If \(A<L\), split at
\(a_*=(A/L)^{1/\beta}\).  The lower part equals \(La_*^\beta=A\);
the upper part is
\[
 \beta A\int_{a_*}^{1}\frac{da}{a}=A\log(L/A).
\]
Thus the integral is at most
\(A[1+\log_+(L/A)]\le A(1+\log L)\), where
\(\log_+x=\max\{\log x,0\}\).  This proves both cases.
\end{proof}

\subsection{The cost of collisions}

The number \(V\) of collisions through the whole cap is fixed under
\(\mathcal E\).  With no collisions all \(k\) root tests are
available.  One collision may remove one test; two or more collisions
will be handled by the deterministic cap.

\begin{proposition}[Conditional squared delay]
\label{prop:conditional-squared-delay}
For \(L\ge1\), set
\[
 h_k(L)=
 \begin{cases}
 1+\log L,&k=3,\\
 1,&k\ge4.
 \end{cases}
\]
Then, almost surely,
\begin{equation}\label{eq:conditional-squared-delay}
 \EE[D_i^2\mid\mathcal E]
 \le C_kd^{2\beta}h_k(L)\mathbf1_{\{V=0\}}
      +C_kLd^\beta\mathbf1_{\{V=1\}}
      +L^2\mathbf1_{\{V\ge2\}}.
\end{equation}
\end{proposition}

\begin{proof}
Because \(W_m\) is \(\mathcal G_m\)-measurable, the tower property
and the root-flip test bound imply, for every integer \(1\le j\le k\)
on the \(\mathcal E\)-measurable event \(\{|J_m|\ge j\}\),
\[
 \EE[\mathbf1_{A_m}W_m\mid\mathcal E]
 \le\EE[\mathbf1_{\{S_m\ne\varnothing\}}W_mq_m^j
         \mid\mathcal E].
\]
Here we used \(q_m\le1\) if more than \(j\) tests are available.

If \(V=0\), take \(j=k\) and use
\(W_mq_m^k\le C_kd^\beta q_m^\gamma\) on nonempty sets, as above.
Substituting in
\eqref{eq:conditional-weighted-delay} and applying integrated
occupation gives \(C_kd^{2\beta}h_k(L)\).

If \(V=1\), at least \(k-1\) tests are available by
\eqref{eq:available-root-tests}.  Since \(k-1\ge\beta\) for \(k\ge3\),
\[
 \mathbf1_{\{S_m\ne\varnothing\}}W_mq_m^{k-1}
 \le C_kd^\beta.
\]
Use this bound at each of the at most \(L\) indices in
\eqref{eq:conditional-weighted-delay}.  It gives \(C_kLd^\beta\).
Finally, when \(V\ge2\), the pathwise bound \(0\le D_i\le L\)
gives \(D_i^2\le L^2\).  This includes environments with no
available test.
\end{proof}

The one-collision contribution contains \(d^\beta\).  Its average
therefore needs the joint distribution of incidences and collisions,
not just an upper bound on the probability of a collision.

\subsection{Joint moments of incidences and collisions}

\begin{lemma}[Weighted root counts]\label{lem:weighted-root-counts}
Fix integers \(k\ge3\) and \(n\ge2k\), a real \(B>0\), and a set
of \(k\) roots.  Take \(M\) independent uniform proper \(k\)-clauses
on \(n\) variables, where \(M\) is a nonnegative integer with
\(M\le Bn\).  Let \(Z\) count clauses
meeting a root, \(V\) count clauses meeting at least two roots, and
\(d=\max\{1,Z\}\).  For every fixed real \(p\ge0\),
\begin{equation}\label{eq:incidence-moments}
 \EE d^p\le C_{k,B,p}.
\end{equation}
For every fixed positive integer \(r\),
\begin{equation}\label{eq:weighted-collision-moments}
 \EE[d^p\mathbf1_{\{V\ge r\}}]\le C_{k,B,p,r}n^{-r}.
\end{equation}
\end{lemma}

\begin{proof}
For one clause, let \(I\) indicate meeting a root and let \(J\)
indicate meeting at least two.  Write \(p_1=\PP(I=1)\) and
\(p_2=\PP(J=1)\).  Union bounds over roots and pairs of roots give
\begin{equation}\label{eq:one-trial-root-counts}
 p_1\le\frac{k^2}{n},\qquad
 p_2\le\binom{k}{2}\frac{k(k-1)}{n(n-1)}
      \le\frac{k^2(k-1)^2}{n^2}.
\end{equation}
Thus \(Z\sim\operatorname{Bin}(M,p_1)\), and
\[
 \EE e^Z=(1+p_1(e-1))^M
 \le\exp(Bk^2(e-1)).
\]
For \(K_p=\sup_{x\ge0}\max\{1,x\}^pe^{-x}<\infty\), we have
\(d^p\le K_pe^Z\).  This proves \eqref{eq:incidence-moments}.

For the collision bound, retain the dependence between \(Z\) and \(V\).
Write \((a)_r=a(a-1)\cdots(a-r+1)\).  For every nonnegative function
\(f\) on the nonnegative integers and \(1\le r\le M\),
\begin{equation}\label{eq:joint-factorial-identity}
 \EE[f(Z)(V)_r]
 =(M)_r p_2^r\,\EE f(r+X),
 \qquad X\sim\operatorname{Bin}(M-r,p_1).
\end{equation}
To prove the identity, expand \((V)_r\) over ordered \(r\)-tuples
of distinct clause indices.  A collision at each chosen index has
probability \(p_2^r\) and forces its incidence indicator to be one.
The remaining \(M-r\) incidence indicators retain their independent
Bernoulli laws with parameter \(p_1\).  Each ordered tuple has the
same expectation, giving the formula.

Since \(\mathbf1_{\{V\ge r\}}\le(V)_r/r!\), use
\(f(x)=\max\{1,x\}^p\) in \eqref{eq:joint-factorial-identity}:
\begin{align*}
 \EE[d^p\mathbf1_{\{V\ge r\}}]
 &\le\frac{(M)_r p_2^r}{r!}\,
          \EE\max\{1,r+X\}^p\\
 &\le\frac{(M)_r p_2^r}{r!}\,
          K_pe^r\exp(Bk^2(e-1))
 \le C_{k,B,p,r}n^{-r}.
\end{align*}
The last inequality uses \(M\le Bn\) and
\eqref{eq:one-trial-root-counts}.  If \(r>M\), the left side is zero.
The same definitions cover \(M=0\) and \(p=0\).
\end{proof}

These moments are taken with the deleted clause fixed and the other
clauses still random.  Once \(\mathcal E\) is fixed, \(Z\) and \(V\)
are known; their rarity returns only when we average over
\(\mathcal E\).

\subsection{Completing the capped upper bound}

\begin{proof}[Proof of the capped upper bound in Theorem~\ref{thm:variance}]
Suppose first that \(n\ge2k\) and \(L\ge1\), and keep \(C_i\) fixed.
The remaining \(L-1\le Bn\) clauses satisfy
Lemma~\ref{lem:weighted-root-counts}.  Averaging
\eqref{eq:conditional-squared-delay} over \(\mathcal E\), and using
the lemma with \(p=2\beta\) in \eqref{eq:incidence-moments} and
\((p,r)=(\beta,1),(0,2)\) in \eqref{eq:weighted-collision-moments},
gives
\begin{align*}
 \EE[D_i^2\mid C_i]
 &\le C_kh_k(L)\EE[d^{2\beta}\mid C_i]
       +C_kL\EE[d^\beta\mathbf1_{\{V=1\}}\mid C_i]
       +L^2\PP(V\ge2\mid C_i)\\
 &\le C_{k,B}h_k(L)+C_{k,B}Ln^{-1}
                         +C_{k,B}L^2n^{-2}\\
 &\le C_{k,B}h_k(L).
\end{align*}
The constants are uniform in \(i\) and in the value of \(C_i\).
Remove this conditioning and use coordinate omission:
\[
 \Var(T)\le\sum_{i=1}^{L}\EE D_i^2
 \le C_{k,B}Lh_k(L).
\]
Since \(L\le Bn\), this is at most
\[
 C_{k,B}n\ell_k(n),\qquad
 \ell_k(n)=
 \begin{cases}
 \log(en),&k=3,\\
 1,&k\ge4.
 \end{cases}
\]
For \(k=3\), we used \(1+\log L\le C_B\log(en)\) when \(L\ge1\).
If \(L=0\), the capped variable is identically zero.  The finitely
many remaining sizes \(k\le n<2k\) are covered by increasing the
constant, since \(0\le T\le\lfloor Bn\rfloor\).  This proves the
upper bound for every fixed positive \(B\) and every \(n\ge k\).
\end{proof}
